\section{Integrasi \textit{Classical Mutual Information} (CMI)}

\subsection{Redefinisi Imbal Hasil Strategis ($\tilde{\mu}_i$)}
Prosedur redefinisi imbal hasil dalam model ini diawali dengan tahap diskritisasi imbal hasil (\textit{return}) untuk mengidentifikasi \textit{state} biner $s \in \{u, d\}$ pada setiap aset. Imbal hasil harian aset $i$ dihitung melalui rasio perubahan harga 
\begin{equation}
    R_{i,t} = \frac{P_{i,t} - P_{i,t-1}}{P_{i,t-1}}
\end{equation}
atau bisa didapat dari log return:
\begin{equation}
    R_{i,t} = \ln \left( \frac{P_{i,t}}{P_{i,t-1}} \right)
\end{equation}
yang kemudian dipetakan ke dalam \textit{state} $u$ (\textit{up}) jika nilai tersebut melampaui ambang batas (\textit{threshold}) nol, atau \textit{state} $d$ (\textit{down}) untuk kondisi sebaliknya. Pada sistem dengan $N=4$ aset, interaksi pasar direpresentasikan melalui 16 kemungkinan \textit{microstate} gabungan $(s_1, s_2, s_3, s_4)$ yang muncul pada setiap satuan waktu $t$. Distribusi frekuensi gabungan kemudian dibangun untuk mengestimasi probabilitas kemunculan setiap \textit{microstate} selama periode observasi $T$, di mana probabilitas gabungan didefinisikan sebagai 
\begin{equation}
    P(s_1, s_2, s_3, s_4) = \frac{\text{count}(s_1, s_2, s_3, s_4)}{T}
\end{equation}

Tahap selanjutnya melibatkan kalkulasi ekspektasi imbal hasil bersyarat (\textit{conditional expected return}) yang menangkap korelasi mikro-struktural antar-aset secara mendalam. Parameter $\bar{R}_i(s_1, s_2, s_3, s_4)$ didefinisikan sebagai nilai ekspektasi imbal hasil aset $i$ yang dikondisikan secara eksklusif pada kemunculan \textit{microstate} gabungan pasar tertentu. Formulasi imbal hasil strategis $\tilde{\mu}_i$ akhirnya disusun sebagai jumlahan terbobot dari ekspektasi bersyarat tersebut terhadap distribusi probabilitas gabungannya. Pendekatan holistik ini menghasilkan parameter imbal hasil yang merefleksikan interdependensi sistemik antar-aset sebagai berikut:
\begin{equation}
    \tilde{\mu}_i = \sum_{s_1, s_2, s_3, s_4 \in \{u,d\}} P(s_1, s_2, s_3, s_4) \times \bar{R}_i(s_1, s_2, s_3, s_4)
\end{equation}
Sebagai ilustrasi untuk aset pertama, ekspansi dari formulasi parameter imbal hasil strategis tersebut dapat dijabarkan melalui seluruh spektrum konfigurasi pasar:
\begin{equation}
    \tilde{\mu}_1 = P(uuuu)\bar{R}_1(uuuu) + P(uuud)\bar{R}_1(uuud) + \dots + P(dddd)\bar{R}_1(dddd)
\end{equation}

\subsection{Informasi Bersama Orde Tinggi ($N=4$)}
Dalam sistem dengan $N=4$ aset ($X_1, X_2, X_3, X_4$), redundansi informasi total diukur melalui \textit{Interaction Information}. Persamaan umumnya mengikuti prinsip inklusi-eksklusi entropi Shannon:
\begin{equation}
    I(X_1:X_2:X_3:X_4) = \sum_{i} H(X_i) - \sum_{i<j} H(X_i, X_j) + \sum_{i<j<k} H(X_i, X_j, X_k) - H(X_1, X_2, X_3, X_4)
\end{equation}

\subsection{Derivasi Rekursif \textit{Interaction Information}}
Derivasi Persamaan (3) dilakukan secara induktif dengan menggunakan definisi rekursif: $I(X_1: \dots :X_n) = I(X_1: \dots :X_{n-1}) - I(X_1: \dots :X_{n-1} | X_n)$.

\paragraph{Langkah 0: Entropi Shannon ($n=1$)}
Informasi tunggal dari sebuah variabel acak $X$ didefinisikan sebagai entropi Shannon $H(X)$:
\begin{equation}
    I(X) = H(X) = - \sum_{x \in \{u,d\}} P(x) \log_2 P(x)
\end{equation}

\paragraph{Langkah 1: Basis Induksi ($n=2$)}
Informasi bersama dua variabel (\textit{Mutual Information}) adalah:
\begin{equation}
\begin{aligned}
    I(X_1:X_2) &= H(X_1) - H(X_1 | X_2) \\
    &= H(X_1) + H(X_2) - H(X_1, X_2) = \sum_{i=1}^2 H(X_i) - \sum_{i<j \le 2} H(X_i, X_j)
\end{aligned}
\end{equation}

\paragraph{Langkah 2: Ekspansi Orde-3 ($n=3$)}
\begin{equation}
\begin{aligned}
    I(X_1:X_2:X_3) &= I(X_1:X_2) - I(X_1:X_2 | X_3) \\
    &= \sum_{i=1}^3 H(X_i) - \sum_{i<j \le 3} H(X_i, X_j) + H(X_1, X_2, X_3)
\end{aligned}
\end{equation}

\paragraph{Langkah 3: Ekspansi Orde-4 ($n=4$)}
\begin{equation}
\begin{aligned}
    I(X_1:X_2:X_3:X_4) &= I(X_1:X_2:X_3) - I(X_1:X_2:X_3 | X_4) \\
    &= \sum_{i=1}^4 H(X_i) - \sum_{i<j \le 4} H(X_i, X_j) + \sum_{i<j<k \le 4} H(X_i, X_j, X_k) - H(X_1, X_2, X_3, X_4)
\end{aligned}
\end{equation}

\subsection{Reduksi Orde melalui Kendala Kardinalitas $K=2$}
Dalam formulasi optimasi portofolio ini, strategi pemilihan aset didefinisikan melalui variabel biner $x_i \in \{0, 1\}$ untuk semesta aset $i \in \{1, 2, 3, 4\}$. Kendala kardinalitas yang membatasi pemilihan tepat dua aset dari empat kandidat dinyatakan sebagai persamaan linear $\sum_{i=1}^4 x_i = 2$. Secara matematis, kendala ini mengimplikasikan bahwa dalam setiap konfigurasi solusi yang valid, tepat terdapat dua variabel yang bernilai satu sementara dua variabel lainnya bernilai nol. Kondisi ini menjadi fondasi krusial dalam menyederhanakan interaksi informasi orde tinggi yang muncul dari integrasi \textit{Interaction Information} ke dalam fungsi objektif.

Integrasi informasi bersama orde-3 dan orde-4 (\textit{Interaction Information}) secara umum menghasilkan fungsi penalti informasional dalam bentuk polinomial derajat tinggi. Fungsi penalti tersebut dapat dinyatakan sebagai kombinasi linear dari interaksi antar-variabel biner sebagai berikut:
\begin{equation}
    P_{informasi} = \sum_{i<j} c_{ij} x_i x_j + \sum_{i<j<k} c_{ijk} x_i x_j x_k + c_{1234} x_1 x_2 x_3 x_4
\end{equation}
di mana $c$ merupakan konstanta bobot informasi yang bersesuaian dengan orde interaksinya. Namun, keberadaan kendala $K=2$ menyebabkan seluruh suku dengan orde lebih tinggi dari dua akan mengalami keruntuhan nilai secara definitif. Evaluasi terhadap suku orde-3 dan orde-4 dilakukan untuk membuktikan reduksi orde ini melalui analisis ruang solusi yang diizinkan oleh kendala kardinalitas.

\paragraph{Pembuktian Suku Orde-3 Runtuh ($x_i x_j x_k = 0$)}
Untuk sembarang kombinasi tiga aset, misalnya $\{x_1, x_2, x_3\}$, suku interaksi $x_1 x_2 x_3$ hanya bernilai non-nol (satu) jika dan hanya jika $x_1 = x_2 = x_3 = 1$. Kondisi ini mensyaratkan bahwa jumlah dari ketiga variabel tersebut adalah $\sum_{i=1}^3 x_i = 3$. Namun, berdasarkan kendala utama $\sum_{i=1}^4 x_i = 2$, dan mengingat sifat non-negatif dari variabel biner ($x_4 \ge 0$), maka nilai maksimum dari $\sum_{i=1}^3 x_i$ dibatasi pada angka dua. Berdasarkan \textit{Pigeonhole Principle} (Prinsip Sarang Merpati), dalam setiap himpunan bagian yang terdiri dari tiga variabel, setidaknya satu variabel pasti bernilai nol, sehingga secara definitif berlaku:
\begin{equation}
    \forall (i,j,k) \in \{1, \dots, 4\} \implies x_i x_j x_k = 0
\end{equation}

\paragraph{Pembuktian Suku Orde-4 Runtuh ($x_1 x_2 x_3 x_4 = 0$)}
Analisis yang serupa berlaku untuk suku interaksi orde-4, di mana hasil kali $x_1 x_2 x_3 x_4$ hanya dapat bernilai satu jika keempat variabel tersebut bernilai satu secara simultan. Hal ini mengimplikasikan total jumlah variabel $\sum_{i=1}^4 x_i = 4$, yang secara eksplisit melanggar kendala kardinalitas $K=2$. Karena jumlah variabel dikunci pada nilai dua, maka dipastikan terdapat tepat dua variabel yang bernilai nol dalam setiap solusi valid. Lenyapnya suku ini secara otomatis menyisakan hanya interaksi biner yang kompatibel dengan struktur Ising \textit{Hamiltonian} sebagai berikut:
\begin{equation}
    x_1 x_2 x_3 x_4 = 0
\end{equation}

Sebagai konsekuensi dari lenyapnya suku-suku orde tinggi tersebut, formulasi penalti informasi tereduksi menjadi bentuk kuadratik murni. Reduksi ini memungkinkan integrasi informasi ke dalam matriks risiko melalui transformasi \textit{Normalized Mutual Information} (NMI). Integrasi ini memerlukan konsistensi dimensional antara satuan informasi dan kovariansi, yang dicapai melalui prosedur skalarisasi. Tinjau dimensi dari entropi Shannon $H(X_i)$ dan \textit{Mutual Information} $I(X_i : X_j)$ yang didefinisikan dalam basis logaritma biner:
\begin{equation}
    H(X_i) = -\sum_{x \in \{u,d\}} P(x) \log_2 P(x) \implies [H] = \text{bit}
\end{equation}
\begin{equation}
    I(X_i : X_j) = H(X_i) + H(X_j) - H(X_i, X_j) \implies [I] = \text{bit}
\end{equation}
Secara teoritis, informasi bersama antara dua variabel acak dibatasi secara simetris oleh rata-rata geometrik dari entropi masing-masing variabel berdasarkan \textit{Upper Bound Theorem}:
\begin{equation}
    I(X_i : X_j) \le \min(H(X_i), H(X_j)) \le \sqrt{H(X_i) H(X_j)}
\end{equation}

Skalarisasi dicapai dengan mendefinisikan NMI sebagai rasio terhadap batas atas teoritisnya, yang secara otomatis membatalkan unit dimensional (\textit{dimension cancellation}):
\begin{equation}
    NMI(i,j) = \frac{I(X_i : X_j)}{\sqrt{H(X_i) H(X_j)}} \implies \left[ \frac{\text{bit}}{\sqrt{\text{bit} \cdot \text{bit}}} \right] = 1
\end{equation}
Hasil akhir dari transformasi ini adalah besaran skalar nirdimensi (\textit{dimensionless}) yang berada pada rentang $NMI(i,j) \in [0, 1]$. Validitas rentang tersebut dibuktikan melalui dua kondisi batas ekstrem:
\begin{itemize}
    \item \textbf{Independensi Statistik ($X_i \perp X_j$):} Karena $I(X_i:X_j) = 0$, maka $\frac{0}{\sqrt{H_i H_j}} = 0$.
    \item \textbf{Identitas Sempurna ($X_i = X_j$):} Karena $I(X_i:X_j) = H(X_i) = H(X_j)$, maka $\frac{H_i}{\sqrt{H_i \cdot H_i}} = 1$.
\end{itemize}

Redefinisi elemen matriks kovariansi akhirnya dirumuskan sebagai penguatan (\textit{amplification}) dari varians tradisional menggunakan skalar NMI tersebut. Formulasi ini menjamin konsistensi matematis saat menggabungkan metrik berbasis informasi dengan metrik berbasis varians-kovarians:
\begin{equation}
    \tilde{\sigma}_{ij} = \sigma_{ij} [1 + NMI(i,j)]
\end{equation}
Melalui skema ini, kontribusi informasional diintegrasikan secara intrinsik ke dalam struktur risiko, sehingga model tetap \textit{robust} terhadap variasi magnitudo data tanpa memerlukan parameter penskalaan eksternal yang bersifat \textit{ad-hoc}.
